\begin{corollary}[Observable projective separation guard]
\label{cor:aesp-cd-projective-separation-guard}
The slope-separation hypothesis in
\cref{cor:aesp-cd-slope-separated-projective-meld} is certifiable without
scanning the child rows.  Store each positive-denominator child range by its
local endpoints $r^-,r^+$ and its lazy matrix $\mathcal P(P,u)$.  Evaluating
\begin{equation*}
 \frac{P_{11}r^-+P_{21}}{P_{12}r^-+P_{22}},\qquad
 \frac{P_{11}r^++P_{21}}{P_{12}r^++P_{22}}
\end{equation*}
and ordering the two values gives the exact pulled-back interval.  Comparing
the two child intervals is therefore an $O(1)$ observable guard.  If they are
disjoint, update the bridge and continue the logarithmic reporter.  If they
overlap, reject this certificate and invoke the declared safe fallback.

Along any trace on which every affected meld passes, the work bound of
\cref{cor:aesp-cd-slope-separated-projective-meld} is implementable rather
than merely promised.  A later common ancestor pullback cannot invalidate a
passed order: a nondegenerate map preserves or reverses both intervals, while
a rank-deficient map collapses their slopes and is handled by retaining the
largest intercept in that common stratum.  The $a_2=0$ endpoint stratum is a
one-dimensional threshold and is maintained separately.
\end{corollary}

This guard does not repair an overlapping meld.  Its value is to make the
successful structural branch observable and safely restartable.
